\section{Related Work}
\textbf{Benchmarking retrieval.}
Existing information retrieval (IR) datasets are typically constructed for information-seeking tasks, such as question-answering~\citep{voorhees2000building_trec,craswell2020overview,kwiatkowski2019natural,chen-etal-2017-reading,maia201818}, claim verification~\citep{thorne2018fever,diggelmann2020climate,wadden2020fact}, or entity retrieval~\citep{hasibi2017dbpedia,petroni-etal-2021-kilt}. 
Recent works expand retrieval benchmarks with more scenarios, such as instruction-following \citep{su-etal-2023-one,weller2024followir,oh2024instructir,li2024url}, multi-hop~\citep{yang2018hotpotqa}, and long-context retrieval~\citep{saad2024benchmarking,zhu2024longembed}. 
Comprehensive benchmarks like BEIR~\citep{thakur2021beir} evaluate retrieval systems on diverse domains and tasks, with relevant documents sharing high semantic overlap with the query. 
Closest to our work, BIRCO~\citep{wang2024birco} is designed to evaluate retrieval systems based on multifaceted objectives by leveraging existing datasets. 
However, it is limited to the LLM reranking setting and uses only a small candidate pool ($\sim 100$ documents) for each query. 
RAR-b~\citep{xiao2024rarb} adapts existing commonsense, math, and code datasets into a retrieval setting to test whether models can directly retrieve answers to reasoning problems. 
However, we focus on a more realistic scenario where the answers are unlikely to be found as a substring in documents.\footnote{We compare to RAR-b in detail in Appendix~\ref{app:rar-b}.}
% However, this setup is somewhat artifiecial, as the answers are typically short strings and do not represent realistic retrieval scenarios where exact matches are unlikely to be found within a document. 
While both of the benchmarks focus on reasoning-intensive retrieval, \ours is the first benchmark to collect realistic user queries and align them with relevant natural documents from large corpora, requiring deep reasoning to identify the correct matches.


% ~\citep{sap-etal-2019-social,zellers-etal-2019-hellaswag,Bisk2019PIQARA,cobbe2021training,hendrycks2021measuring,muennighoff2023octopack} 
\textbf{Dense retrieval models and retrieval augmented generation.} State-of-the-art retrieval systems often use dense models to encode text with rich representation. 
These models are trained on unsupervised data \citep{lee-etal-2019-latent,izacard2022unsupervised}, supervised data \citep{su-etal-2023-one,asai-etal-2023-task,muennighoff2022sgpt}, as well as LLM-generated data~\citep{lee2024gecko,wang2023improving,muennighoff2024generative}. 
In this work, we benchmark a diverse set of models across different axes: sparse and dense; small and large; open-source and proprietary. 
Additionally, as dense generative models continue to improve, retrieval-augmented generation ~\citep[RAG;][]{asai2020reliable, borgeaud2022retro, asai2023self, muennighoff2024generative,asai2024reliable,gao2023enabling,shi2024replug}, which retrieves relevant documents to help generate coherent answers, has become an important application. 
% In this work, we focus on retrieval and leave the exploration of RAG evaluations on \ours{} for future work. 
% We conduct initial analyses and demonstrate in Appendix~\ref{sec:case_study} how retrieving relevant documents can help improve model generation for reasoning-intensive tasks.
Although we mainly focus on retrieval in this work, we conduct initial analyses and demonstrate that using stronger retrievers improves model generation for reasoning-intensive tasks.
% We conduct initial analyses and demonstrate in Appendix~\ref{sec:case_study} how retrieving relevant documents can help improve model generation for reasoning-intensive tasks.

%\vspace{-2mm}
\textbf{Benchmarking reasoning.} 
Many benchmarks aim to evaluate the reasoning abilities of LLMs, especially focused on mathematics and coding. As for mathematics, for example, datasets include GSM8K~\citep{cobbe2021training} and its extensions GSM1K~\citep{zhang2024careful}, TheoremQA~\citep{chen-etal-2023-theoremqa}, MATH~\citep{hendrycks2021measuring}, and LeanDojo~\citep{yang2023leandojo}. 
As for coding, %benchmarks such as
HumanEval~\citep{chen2021evaluating}, MBPP~\citep{austin2021program}, and LiveCodeBench~\citep{jain2024livecodebench} are often used. 
These benchmarks contain question-answer pairs and are usually sourced from textbooks, online resources, competitions, or domain experts. 
We source queries from selected high-quality datasets and construct \ours{} through additional annotations, creating a realistic reasoning-intensive retrieval benchmark.
